% =============================================================================
% RAPPORT PROFESSIONNEL — UrbanOps (COMPLET — V2.0)
% Compilation : pdflatex x2 (ou xelatex x2) depuis docs/latex/
% =============================================================================
\documentclass[12pt,a4paper,twoside,openright]{book}

% --- Encodage & langue ---
\usepackage[T1]{fontenc}
\usepackage[utf8]{inputenc}
\usepackage[french]{babel}
\usepackage{microtype}

% --- Mise en page & texte ---
\usepackage[a4paper,left=2.5cm,right=2.5cm,top=2.5cm,bottom=2.5cm,bindingoffset=0.5cm]{geometry}
\usepackage{setspace}
\setstretch{1.22}
\usepackage{ragged2e}
\usepackage{enumitem}
\setlist{itemsep=0.35em, parsep=0pt, topsep=0.4em}

% --- Graphiques & couleurs ---
\usepackage{graphicx}
\usepackage[table]{xcolor}
\definecolor{UrbanBlue}{HTML}{1E3A5F}
\definecolor{UrbanAccent}{HTML}{2563EB}
\definecolor{UrbanGray}{HTML}{4B5563}
\definecolor{LightBand}{HTML}{F3F4F6}
\definecolor{CodeBg}{HTML}{F8FAFC}
\definecolor{AlertRed}{HTML}{DC2626}
\definecolor{SuccessGreen}{HTML}{16A34A}

\usepackage{float}
\usepackage{caption}
\captionsetup{font=small,labelfont={bf,color=UrbanBlue}}

% --- Tableaux ---
\usepackage{booktabs}
\usepackage{array}
\usepackage{longtable}
\usepackage{tabularx}
\newcolumntype{Y}{>{\RaggedRight\arraybackslash}X}

% --- Maths ---
\usepackage{amsmath,amssymb}

% --- Code & encadrés ---
\usepackage{listings}
\lstdefinestyle{urban}{%
  backgroundcolor=\color{CodeBg},
  basicstyle=\ttfamily\footnotesize,
  breaklines=true,
  frame=single,
  framerule=0.4pt,
  rulecolor=\color{UrbanBlue!40},
  columns=fullflexible,
  keepspaces=true,
  showstringspaces=false,
  tabsize=2,
}
\lstset{style=urban}
\lstdefinestyle{sqlurban}{style=urban,language=SQL,keywordstyle=\color{UrbanAccent},stringstyle=\color{AlertRed}}
\lstdefinestyle{javaurban}{style=urban,language=Java,keywordstyle=\color{UrbanAccent},stringstyle=\color{AlertRed}}
\lstdefinestyle{bashurban}{style=urban,language=bash}

\usepackage[most]{tcolorbox}
\tcbuselibrary{skins,breakable}

% --- Apparence table des matières ---
\usepackage{tocloft}
\renewcommand{\cftchapfont}{\bfseries\color{UrbanBlue}}
\renewcommand{\cftsecfont}{\color{UrbanGray}}
\setlength{\cftbeforechapskip}{0.55em}

% --- Titres (numérotés = arabe 1, 2… après introduction) ---
\usepackage{titlesec}
\titleformat{\chapter}[display]
  {\normalfont\huge\bfseries\color{UrbanBlue}}
  {\chaptertitlename~\thechapter}
  {12pt}
  {\Huge}
\titlespacing*{\chapter}{0pt}{-8pt}{22pt}

\titleformat{name=\chapter,numberless}[display]
  {\normalfont\huge\bfseries\color{UrbanBlue}}
  {}
  {0pt}
  {\Huge}
\titlespacing*{name=\chapter,numberless}{0pt}{-8pt}{22pt}

\titleformat{\section}
  {\Large\bfseries\color{UrbanAccent}}
  {\thesection}{0.75em}{}
\titlespacing*{\section}{0pt}{1.1em}{0.65em}

\titleformat{\subsection}
  {\large\bfseries\color{UrbanGray}}
  {\thesubsection}{0.65em}{}

\setcounter{tocdepth}{3}
\setcounter{secnumdepth}{3}

% --- Liens PDF ---
\usepackage[hidelinks,pdfusetitle]{hyperref}
\hypersetup{
  colorlinks=true,
  linkcolor=UrbanBlue,
  urlcolor=UrbanAccent,
  citecolor=UrbanBlue,
  bookmarksnumbered=true,
  pdftitle={UrbanOps - Rapport Technique Complet},
  pdfauthor={Equipe UrbanOps},
}
\usepackage{bookmark}
\usepackage[nameinlink]{cleveref}

% --- En-têtes / pieds ---
\usepackage{fancyhdr}
\pagestyle{fancy}
\fancyhf{}
\fancyhead[LE]{\small\textcolor{UrbanGray}{UrbanOps — Supervision urbaine}}
\fancyhead[RO]{\small\textcolor{UrbanGray}{2025–2026}}
\fancyhead[LO,RE]{\small\leftmark}
\fancyfoot[C]{\thepage}
\renewcommand{\headrulewidth}{0.45pt}
\renewcommand{\footrulewidth}{0pt}
\makeatletter
\let\ps@plain\ps@fancy
\makeatother

% =============================================================================
% Encadrés réutilisables — SANS BACKGROUND, contours uniquement
% =============================================================================
\newtcolorbox{resumebox}{
  enhanced,
  colback=white,
  colframe=UrbanBlue,
  boxrule=0.6pt,
  arc=3mm,
  left=8mm,right=8mm,top=6mm,bottom=6mm,
  breakable,
}

\newtcolorbox{importantbox}{
  enhanced,
  colback=white,
  colframe=AlertRed,
  boxrule=0.6pt,
  arc=3mm,
  left=8mm,right=8mm,top=6mm,bottom=6mm,
  breakable,
}

\newtcolorbox{successbox}{
  enhanced,
  colback=white,
  colframe=SuccessGreen,
  boxrule=0.6pt,
  arc=3mm,
  left=8mm,right=8mm,top=6mm,bottom=6mm,
  breakable,
}

\newtcolorbox{archibox}[1][]{
  enhanced,
  colback=white,
  colframe=UrbanBlue,
  fonttitle=\bfseries,
  title={#1},
  attach boxed title to top center={yshift=-2mm},
  boxed title style={colback=UrbanBlue,colframe=UrbanBlue},
  arc=2mm,
  breakable,
}

% =============================================================================
% Figures : inclusion ou emplacement réservé
% =============================================================================
\newcommand{\urbanfig}[4]{%
  \begin{figure}[htbp]%
    \centering%
    \IfFileExists{#2}{%
      \includegraphics[width=#3\textwidth]{#2}%
    }{%
      \begin{tcolorbox}[
        enhanced,colback=white,colframe=UrbanBlue,
        width=0.92\textwidth,arc=2mm,halign=center,valign=center,
        height=0.30\textheight,boxrule=0.6pt]
        {\large\bfseries Emplacement figure}\\[3mm]
        {\small\texttt{#2}}\\[2mm]
        {\itshape Insérer l'image exportée (PNG/PDF).}\\[1mm]
        {\footnotesize Légende : \textbf{#4}}
      \end{tcolorbox}%
    }%
    \caption{#4}%
    \label{#1}%
  \end{figure}%
}

% --- Chapitres I à IV (romains) ---
\newcounter{prechapter}
\setcounter{prechapter}{0}
\newcommand{\ChapterRoman}[1]{%
  \refstepcounter{prechapter}%
  \phantomsection
  \chapter*{Chapitre \Roman{prechapter} — #1}%
  \addcontentsline{toc}{chapter}{Chapitre \Roman{prechapter} — #1}%
  \markboth{Chapitre \Roman{prechapter} — #1}{}%
}

\newcommand{\SecRoman}[1]{%
  \phantomsection
  \section*{#1}%
  \addcontentsline{toc}{section}{#1}%
}

\title{Rapport technique — UrbanOps\\\large Plateforme de supervision urbaine intelligente}
\author{Équipe UrbanOps}
\date{\today}

\begin{document}

% =============================================================================
% PAGE DE GARDE
% =============================================================================
\begin{titlepage}
  \thispagestyle{empty}
  \newgeometry{margin=2.2cm}
  \centering
  \null\vspace*{0.8cm}

  {\color{AlertRed}\rule{0.82\textwidth}{2.2pt}}\par\vspace{0.75cm}

  {\Large\bfseries\textsc{Université / École}}\par\vspace{0.2cm}
  {\large Filière Génie Logiciel / Informatique}\par\vspace{0.35cm}
  {\normalsize\itshape Rapport de Projet Final — PFE 2025–2026}\par\vspace{1.6cm}

  {\Huge\bfseries\color{UrbanBlue} UrbanOps}\par\vspace{0.45cm}
  {\LARGE\bfseries Plateforme intelligente de supervision urbaine}\par\vspace{0.25cm}
  {\large Signalements géolocalisés, alertes ciblées, IA et tableau de bord}\par\vspace{1.4cm}

  {\color{AlertRed}\rule{0.55\textwidth}{0.55pt}}\par\vspace{0.65cm}
  {\LARGE\bfseries Rapport technique complet}\par\vspace{0.55cm}
  {\color{AlertRed}\rule{0.55\textwidth}{0.55pt}}\par\vspace{1.6cm}

  \begin{minipage}{0.88\textwidth}\centering
    \color{UrbanGray}\large
    Architecture empilée \textbf{3~tiers} (Next.js~14, Spring~Boot~3, PostgreSQL), sécurité \textbf{JWT}, module d'\textbf{IA} (Gemini + repli), \textbf{4~middlewares} (REST, JMS, RMI, SOAP), tests \textbf{JUnit~5}, couverture \textbf{JaCoCo}, analyse \textbf{SonarQube}.
  \end{minipage}\par\vfill

  {\color{UrbanBlue}\rule{0.82\textwidth}{0.6pt}}\par\vspace{0.65cm}

  {\centering
  \renewcommand{\arraystretch}{1.35}
  \begin{tabular}{@{}>{\bfseries}r @{\hspace{1em}} p{0.52\textwidth}@{}}
    Auteur(s) :     & Équipe de développement UrbanOps \\
    Encadrant(e) :  & Professeur [À remplir] \\
    Technologies :  & Spring Boot~3.2, Next.js~14, PostgreSQL, Gemini, JWT \\
    Année :         & 2025 – 2026 \\
  \end{tabular}\par}

  \vspace{0.9cm}
  {\large\textbf{Version}~2.0 \quad\textemdash\quad \today}\par
  \vspace*{0.6cm}
  \restoregeometry
\end{titlepage}

\cleardoublepage

% =============================================================================
% PRÉLIMINAIRES
% =============================================================================
\frontmatter

% --- Dédicaces ---
\chapter*{Dédicaces}
\addcontentsline{toc}{chapter}{Dédicaces}
\thispagestyle{empty}
\vspace*{2.8cm}
\begin{center}
\begin{minipage}{0.82\textwidth}\centering\itshape\large
À tous les contributeurs de villes plus intelligentes et plus résilientes.\\[1.2em]
Aux citoyens et autorités qui inspireront cette plateforme.\\[1.5em]
\ldots
\end{minipage}
\end{center}
\cleardoublepage

% --- Remerciements ---
\chapter*{Remerciements}
\addcontentsline{toc}{chapter}{Remerciements}
\begin{center}
\begin{minipage}{0.88\textwidth}
\justifying
Nous exprimons notre profonde gratitude à notre encadrant(e) pour les relectures minutieuses, les conseils méthodologiques avisés et la disponibilité constante tout au long du projet.

Nous remercions les membres du jury pour l'attention bienveillante portée à ce rapport et pour leurs évaluations constructives.

Les excellentes communautés open-source autour de \textbf{Spring Framework}, \textbf{Next.js}, \textbf{PostgreSQL} et \textbf{Google Gemini} ont mis à disposition une documentation riche et des outils essentiels à la réalisation de ce travail.

Merci également à nos collègues pour les discussions et les retours tout au long du développement.
\end{minipage}
\end{center}
\cleardoublepage

% --- Résumé ---
\chapter*{Résumé}
\addcontentsline{toc}{chapter}{Résumé}
\thispagestyle{plain}
\begin{resumebox}
\justifying
\textbf{UrbanOps} est une plateforme \textbf{full-stack} dédiée au signalement, au suivi et à la supervision d'incidents urbains (voirie dégradée, distribution d'eau interrompue, éclairage déficient, accumulation de déchets, problèmes de sécurité, etc.).

L'architecture repose sur un backend \textbf{Spring Boot~3} exposant une API REST versionnée (\texttt{/api/v1}), sécurisée par jetons \textbf{JWT}, avec persistance fiable en \textbf{PostgreSQL} via \textbf{JPA/Hibernate}. Un module d'analyse textuelle s'appuie sur \textbf{Google Gemini} pour classifier les incidents, avec mécanisme de \textbf{repli par règles} en cas d'indisponibilité. Le frontend \textbf{Next.js~14} (TypeScript, Tailwind, Axios) offre une landing page, un tableau de bord interactif, une carte \textbf{Leaflet} et des formulaires réactifs.

Le système intègre \textbf{quatre middlewares} : \textbf{REST} (API HTTP/JSON), \textbf{JMS} (messages asynchrones ActiveMQ), \textbf{RMI} (procédures distantes Java) et \textbf{SOAP} (services web XML). Les tests couvrent \textbf{JUnit~5} (mocks Mockito), avec mesure de couverture \textbf{JaCoCo} et analyse qualité \textbf{SonarQube}. La journalisation, la sécurité et le monitoring via Actuator complètent le système.

Ce rapport détaille besoins, conception \textbf{UML}, implémentation, middlewares, tests et pilotage qualité.
\end{resumebox}

\vspace{0.9cm}
\begin{center}
\textbf{Mots-clés :} supervision urbaine ; ville intelligente ; Spring Boot ; Next.js ; PostgreSQL ; JWT ; IA ; Gemini ; REST ; JMS ; RMI ; SOAP ; JUnit~5 ; JaCoCo ; SonarQube.
\end{center}
\cleardoublepage

% --- Abstract ---
\chapter*{Abstract}
\addcontentsline{toc}{chapter}{Abstract}
\thispagestyle{plain}
\begin{center}
\begin{minipage}{0.88\textwidth}
\justifying
\textit{UrbanOps is a comprehensive full-stack smart-city platform for incident reporting, monitoring, and supervision in an urban context. The backend architecture integrates Spring Boot~3 with REST APIs, JWT authentication, and PostgreSQL persistence via JPA. AI-assisted triage leverages the Google Gemini API with deterministic fallback rules. The Next.js~14 frontend provides responsive dashboards and interactive maps. Four middleware layers are implemented: REST (HTTP/JSON), JMS (asynchronous messaging), RMI (remote method invocation), and SOAP (XML web services). Quality assurance includes JUnit~5 testing, JaCoCo coverage analysis, and SonarQube static analysis. This document presents requirements, UML design, implementation details, middleware integration, testing strategies, and quality gates.}
\end{minipage}
\end{center}

\vspace{0.9cm}
\begin{center}
\textbf{Keywords:} urban supervision; smart city; Spring Boot; Next.js; PostgreSQL; JWT; AI; Gemini; REST; JMS; RMI; SOAP; JUnit~5; JaCoCo; SonarQube.
\end{center}
\cleardoublepage

% --- Abréviations ---
\chapter*{Liste des abréviations et acronymes}
\addcontentsline{toc}{chapter}{Liste des abréviations et acronymes}
\begin{table}[H]
  \centering
  \renewcommand{\arraystretch}{1.18}
  \begin{tabularx}{0.96\textwidth}{@{}l Y@{}}
    \toprule
    \textbf{Sigle} & \textbf{Définition} \\
    \midrule
    API & Interface de programmation d'application \\
    AMQP & Advanced Message Queuing Protocol \\
    AI / IA & Intelligence artificielle \\
    BCrypt & Fonction de hachage de mots de passe \\
    DAO & Objet d'accès aux données \\
    DTO & Objet de transfert de données \\
    FTP & Protocole de transfert de fichiers \\
    GANTT & Diagramme de planification par tâches \\
    GPS & Système mondial de positionnement \\
    HTTP/HTTPS & Protocole(s) de transfert hypertexte \\
    IDE & Environnement de développement intégré \\
    JaCoCo & Java Code Coverage \\
    JMS & Java Message Service \\
    JPA & Jakarta Persistence API (ORM) \\
    JSON & Notation d'objet JavaScript \\
    JWT & JSON Web Token \\
    ORM & Mapping objet-relationnel \\
    PERT & Program Evaluation and Review Technique \\
    REST & Transfert d'état représentatif \\
    RMI & Remote Method Invocation \\
    SMTP & Protocole d'envoi de courriels \\
    SOAP & Architecture orientée services \\
    SQL & Langage de requête structuré \\
    TLS & Couche de transport sécurisée \\
    UML & Langage de modélisation unifié \\
    XML & Langage extensible de balisage \\
    \bottomrule
  \end{tabularx}
\end{table}
\cleardoublepage

\listoffigures
\cleardoublepage

\listoftables
\cleardoublepage

\tableofcontents
\cleardoublepage

% =============================================================================
% CHAPITRES I À IV (PRÉAMBULE, numérotation ROMAINE)
% =============================================================================
\mainmatter

\ChapterRoman{Cadre général et problématique}
\label{chap:I}
\begin{center}
\begin{minipage}{0.92\textwidth}
\justifying
Ce chapitre pose le \textbf{contexte métropolitain} (Marrakech, Maroc), la \textbf{problématique} (signalements éparpillés, délais de traitement imprévisibles, priorisation absente) et les \textbf{objectifs} d'UrbanOps : centraliser, tracer, alerter et offrir une interface citoyenne ergonomique.
\end{minipage}
\end{center}

\SecRoman{Contexte et enjeux}
La transformation numérique des services urbains impose des canaux \textbf{unifiés} pour la remontée d'information et une \textbf{traçabilité} exhaustive des traitements. Les autorités et citoyens nécessitent une \textbf{visibilité en temps réel} et une \textbf{responsabilité partagée}.

\SecRoman{Objectifs du projet}
\begin{itemize}
  \item Permettre la \textbf{déclaration} d'incidents géolocalisés avec médias (photos, vidéos) ;
  \item Offrir un \textbf{tableau de bord} centralisé et des \textbf{alertes ciblées} vers les autorités ;
  \item Garantir une \textbf{piste d'audit} immuable (\texttt{INC-xxxx}, statuts, horodatage UTC) ;
  \item Intégrer \textbf{l'IA} pour classifier automatiquement les incidents.
\end{itemize}

\SecRoman{Périmètre fonctionnel}
Le périmètre couvre le cycle de vie complet d'un incident (création, suivi, résolution), l'authentification par rôles (\texttt{ADMIN}, \texttt{CITIZEN}), les statistiques agrégées et les alertes. \textbf{Hors périmètre} : application mobile native (perspective future), notifications push, intégration IoT.

\SecRoman{Acteurs et cas d'usage} Citoyens (reporters), autorités (gestionnaires), administrateurs (superviseurs), système (IA). Chaque rôle dispose de droits d'accès et de vues différenciées.

\ChapterRoman{Analyse des besoins et spécifications}
\label{chap:II}

\SecRoman{Fonctionnalités principales}
\begin{longtable}{@{}p{3.4cm} >{\RaggedRight\arraybackslash}p{10.8cm}@{}}
  \toprule
  \textbf{Fonctionnalité} & \textbf{Description} \\
  \midrule
  \endfirsthead
  \toprule
  \textbf{Fonctionnalité} & \textbf{Description} \\
  \midrule
  \endhead
  Authentification & Inscription, connexion, génération JWT, récupération profil (\texttt{/auth/me}). \\
  Signalement & Création multipart (\texttt{POST /incidents}) : JSON + fichier photo, géolocalisation. \\
  Classification IA & Analyse Gemini pour sévérité (\texttt{HIGH}, \texttt{MEDIUM}, \texttt{LOW}). \\
  Carte interactive & Projection légère \texttt{GET /incidents/map} avec Leaflet. \\
  Tableau de bord & Agrégats \texttt{GET /stats/dashboard} : comptages, tendances. \\
  Alertes & Création et envoi e-mail automatis. pour sévérités \texttt{HIGH} / \texttt{MEDIUM}. \\
  Administration & Gestion utilisateurs, statut incident, alertes (rôle \texttt{ADMIN}). \\
  Endpoints multiples & Filtrages avancés, pagination, tri. \\
  \bottomrule
\end{longtable}

\SecRoman{Besoins non fonctionnels}
\begin{itemize}
  \item \textbf{Performance} : pagination, projections légères pour la carte, caches Redis (optionnel) ;
  \item \textbf{Sécurité} : mots de passe \texttt{BCrypt}, filtre JWT, règles Spring Security, HTTPS/TLS ;
  \item \textbf{Maintenabilité} : couches controller / service / repository, documentation OpenAPI ;
  \item \textbf{Disponibilité} : Actuator \texttt{/health}, journalisation fichier, repli IA déterministe ;
  \item \textbf{Observabilité} : logs structurés, traces des opérations critiques.
\end{itemize}

\SecRoman{Règles métier}
\begin{itemize}
  \item Sévérités : \texttt{HIGH} (rouge, alerte immédiate), \texttt{MEDIUM} (orange, alerte différée), \texttt{LOW} (vert, consultation) ;
  \item Statuts : \texttt{OPEN}, \texttt{IN\_PROGRESS}, \texttt{RESOLVED}, \texttt{REJECTED} ;
  \item Rôles : \texttt{ADMIN} (tous les droits), \texttt{CITIZEN} (lecture propres incidents).
\end{itemize}

\ChapterRoman{Conception et architecture}
\label{chap:III}

\SecRoman{Architecture tres-tiers}
\begin{archibox}[Vue logique empilée]
\begin{center}
\renewcommand{\arraystretch}{1.25}
\begin{tabular}{|>{\centering\arraybackslash}p{0.82\textwidth}|}
\hline
\textbf{Frontend} — Next.js~14 (TypeScript, Tailwind, Axios, Leaflet) \\
\hline
\textbf{API REST} — Spring~MVC sur \texttt{/api/v1} \\
\hline
\textbf{Métier} — Services transactionnels, règles, IA, mails \\
\hline
\textbf{Accès données} — Spring~Data~JPA avec Specifications \\
\hline
\textbf{PostgreSQL} — Schéma généré/mis à jour par Hibernate \\
\hline
\end{tabular}
\end{center}
\end{archibox}

\SecRoman{Entités JPA principales}
Cinq entités centrales : \texttt{User}, \texttt{Category}, \texttt{Sector}, \texttt{Incident}, \texttt{Alert}. Des indices sur les colonnes \texttt{status}, \texttt{severity}, \texttt{latitude}/\texttt{longitude} optimisent les requêtes.

\SecRoman{Diagrammes UML à insérer}
\begin{itemize}
  \item \texttt{figures/gantt\_planning.png} — Diagramme GANTT (planification du projet, jalons) ;
  \item \texttt{figures/pert\_taches.png} — Diagramme PERT (dépendances, chemins critiques) ;
  \item \texttt{figures/uml\_classes\_metier.png} — Diagramme de classes (domaine métier) ;
  \item \texttt{figures/uml\_classes\_technique.png} — Diagramme de classes (couches techniques) ;
  \item \texttt{figures/uml\_cas\_usage.png} — Diagramme de cas d'utilisation ;
  \item \texttt{figures/uml\_composants.png} — Diagramme par composants ;
  \item \texttt{figures/uml\_sequence\_incident.png} — Diagramme de séquence (création incident).
\end{itemize}

\urbanfig{fig:gantt1}{figures/gantt_planning.png}{0.95}{Diagramme de Gantt — jalons et phases du projet.}
\urbanfig{fig:pert1}{figures/pert_taches.png}{0.92}{Diagramme PERT — tâches, dépendances et cheminement critique.}
\urbanfig{fig:classmetier1}{figures/uml_classes_metier.png}{0.95}{Diagramme de classes — modèle métier (User, Incident, Alert, etc.).}
\urbanfig{fig:classtech1}{figures/uml_classes_technique.png}{0.95}{Diagramme de classes — couches technique (Controller, Service, Repository).}
\urbanfig{fig:usecase1}{figures/uml_cas_usage.png}{0.88}{Diagramme de cas d'utilisation — interacteurs et scénarios.}
\urbanfig{fig:compo1}{figures/uml_composants.png}{0.92}{Diagramme par composants — modules, dépendances, interfaces.}
\urbanfig{fig:seq1}{figures/uml_sequence_incident.png}{0.92}{Diagramme de séquence — création incident avec analyse IA et alerte.}

\ChapterRoman{Environnement technique et choix technologiques}
\label{chap:IV}

\SecRoman{Stack backend}
\begin{itemize}
  \item \textbf{Langage} : Java~17 (LTS) ;
  \item \textbf{Framework} : Spring~Boot~3.2.x (Web, Data~JPA, Security, Validation, Mail, Actuator) ;
  \item \textbf{BD} : PostgreSQL~15+ ;
  \item \textbf{ORM} : Hibernate + JPA ;
  \item \textbf{Build} : Maven~3.8+, gestion des dépendances ;
  \item \textbf{Sécurité} : JJWT (JWT), BCrypt, Spring~Security ;
  \item \textbf{IA} : Google Gemini API + repli déterministe ;
  \item \textbf{Email} : Intégration SMTP (Gmail) ;
  \item \textbf{Documentation} : Springdoc~OpenAPI, Swagger~UI ;
  \item \textbf{Utils} : Lombok, MapStruct (optionnel).
\end{itemize}

\SecRoman{Stack frontend}
\begin{itemize}
  \item \textbf{Framework} : Next.js~14 (App~Router) ;
  \item \textbf{Langage} : TypeScript~5 ;
  \item \textbf{Styling} : Tailwind~CSS~3 ;
  \item \textbf{HTTP} : Axios + intercepteurs JWT ;
  \item \textbf{Cartographie} : Leaflet + OpenStreetMap ;
  \item \textbf{Graphiques} : Recharts ;
  \item \textbf{Build} : npm / yarn / bun.
\end{itemize}

\SecRoman{Middleware et protocoles}
\begin{table}[H]
  \centering
  \caption{Les quatre middlewares intégrés dans UrbanOps.}
  \label{tab:middleware}
  \renewcommand{\arraystretch}{1.2}
  \begin{tabularx}{0.95\textwidth}{@{}l l c X@{}}
    \toprule
    \textbf{Middleware} & \textbf{Protocole} & \textbf{Port} & \textbf{Rôle} \\
    \midrule
    REST & HTTP/JSON & 8080 & API CRUD standard, endpoints citoyen/admin \\
    JMS & AMQP (ActiveMQ) & embedded & Messages asynchrones, alertes découplées \\
    RMI & Java RMI & 1099 & Appels distants Java, services métier (IA) \\
    SOAP & HTTP/XML & 8080 & Services web, compatibilité avec systèmes legacy \\
    \bottomrule
  \end{tabularx}
\end{table}

\begin{itemize}
  \item \textbf{REST} : Endpoint principal, stateless, scalable. Endpoints disponibles sur \texttt{http://localhost:8080/api/v1/swagger-ui.html} ;
  \item \textbf{JMS} : Envoi d'alertes mails de manière asynchrone (ActiveMQ embedded) ;
  \item \textbf{RMI} : Service \texttt{AIAnalysisService} exporté sur \texttt{rmi://localhost:1099} ;
  \item \textbf{SOAP} : WSDL exposé à \texttt{/api/v1/soap/urbanops.wsdl} pour intégrations tierces.
\end{itemize}

\SecRoman{Qualité et industrialisation}
\begin{itemize}
  \item \textbf{Tests unitaires} : JUnit~5 (\texttt{@Test}, \texttt{@Disabled}), Mockito (mocks de repositories/services) ;
  \item \textbf{Couverture} : JaCoCo avec \textbf{seuils de qualité} (ex. \texttt{pom.xml} : $\geq 70\%$ couverture ciblée) ;
  \item \textbf{Analyse statique} : SonarQube (bugs, vulnérabilités, détection de code dupliqué) ;
  \item \textbf{Journalisation} : Fichiers de log structurés dans \texttt{backend/logs/} ;
  \item \textbf{Actuator} : Healthchecks, métriques, endpoints de diagnostique.
\end{itemize}

% =============================================================================
% CHAPITRES EN NUMÉROTATION ARABE (1, 2, 3, …)
% =============================================================================
\setcounter{chapter}{0}

\chapter{Introduction générale}
\label{chap:intro}

\section{Problématique et contexte}
Comment offrir une plateforme \textbf{fiable}, \textbf{sécurisée}, \textbf{évolutive} et \textbf{inclusive} pour agréger les signalements citoyens d'incidents urbains et accélérer la réponse opérationnelle des autorités ?

\section{Méthodologie adoptée}
\begin{enumerate}
  \item Analyse approfondie des besoins fonctionnels et non-fonctionnels ;
  \item Conception UML (classes, séquences, cas d'usage, composants) ;
  \item Implémentation incrémentale par les couches (backend → frontend → tests) ;
  \item Tests automatisés JUnit~5 et couverture JaCoCo ;
  \item Analyse statique SonarQube pour la qualité ;
  \item Déploiement et documentation OpenAPI.
\end{enumerate}

\section{Plan du rapport}
Ce document suit la structure suivante :
\begin{itemize}
  \item \textbf{Chapitres I–IV} : Préambule (cadre, besoins, conception, technologies) ;
  \item \textbf{Chapitre 1} : Introduction (présent) ;
  \item \textbf{Chapitre 2} : Réalisation backend et persistance ;
  \item \textbf{Chapitre 3} : Frontend et intégration HTTP ;
  \item \textbf{Chapitre 4} : Sécurité, IA et middlewares ;
  \item \textbf{Chapitre 5} : Tests, couverture JaCoCo, SonarQube ;
  \item \textbf{Chapitre 6} : Déploiement et exploitation ;
  \item \textbf{Chapitre 7} : Conclusion et perspectives.
\end{itemize}

\chapter{Réalisation — Backend et persistance}
\label{chap:real}

\section{Architecture couches du backend}

\subsection{Couche contrôleur (REST)}
Les contrôleurs exposent les ressources via des endpoints REST :
\begin{itemize}
  \item \texttt{AuthController} — \texttt{POST /auth/login}, \texttt{POST /auth/register}, \texttt{GET /auth/me} ;
  \item \texttt{IncidentController} — CRUD incidents, carte légère, filtres avancés ;
  \item \texttt{StatsController} — Agrégats tableaux de bord ;
  \item \texttt{AlertController} — Gestion des alertes (ADMIN) ;
  \item \texttt{UserController} — Gestion utilisateurs ;
  \item \texttt{CategoryController} — Listing catégories ;
  \item \texttt{SectorController} — Listing secteurs.
\end{itemize}

\subsection{Couche service}
Les services encapsulent la logique métier :
\begin{itemize}
  \item \texttt{IncidentService.createIncident} : validation catégorie/secteur, appel IA (Gemini), persistance, alerte ;
  \item \texttt{AIAnalysisService} : appel API Gemini, fallback mots-clés ;
  \item \texttt{AlertService} : création alerts, envoi via \texttt{EmailService} (async) ;
  \item \texttt{AuthService} : génération JWT, validation passwords ;
  \item \texttt{UserService}, \texttt{CategoryService}, \texttt{SectorService} : CRUD métier.
\end{itemize}

\subsection{Couche repository}
Spring~Data~JPA avec projections et Specifications :
\begin{itemize}
  \item \texttt{IncidentRepository} — findAll, findAllForMapProjection (interface projection), Specifications pour filtres ;
  \item \texttt{UserRepository} — findByUsername, findByEmail ;
  \item \texttt{AlertRepository}, \texttt{CategoryRepository}, \texttt{SectorRepository} — CRUD standard.
\end{itemize}

\subsection{Persistance — PostgreSQL}
\begin{lstlisting}[style=sqlurban,caption=Schéma conceptuel incidents.,label=lst:schema]
CREATE TABLE incidents (
  id BIGSERIAL PRIMARY KEY,
  reference_code VARCHAR(20) UNIQUE,
  title VARCHAR(255) NOT NULL,
  description TEXT,
  severity VARCHAR(20),  -- HIGH, MEDIUM, LOW
  status VARCHAR(20),    -- OPEN, IN_PROGRESS, RESOLVED
  latitude DOUBLE PRECISION NOT NULL,
  longitude DOUBLE PRECISION NOT NULL,
  photo_url VARCHAR(255),
  category_id BIGINT NOT NULL REFERENCES categories(id),
  sector_id BIGINT NOT NULL REFERENCES sectors(id),
  reported_by_id BIGINT REFERENCES users(id),
  ai_analysis_result TEXT,
  authority_notified BOOLEAN DEFAULT FALSE,
  created_at TIMESTAMP DEFAULT CURRENT_TIMESTAMP,
  updated_at TIMESTAMP DEFAULT CURRENT_TIMESTAMP,
  resolved_at TIMESTAMP
);

CREATE INDEX idx_incidents_status ON incidents(status);
CREATE INDEX idx_incidents_severity ON incidents(severity);
CREATE INDEX idx_incidents_coordinates ON incidents(latitude, longitude);
\end{lstlisting}

\section{Endpoints REST représentatifs}

\begin{table}[H]
  \centering
  \caption{Endpoints REST principaux d'UrbanOps.}
  \label{tab:endpoints}
  \renewcommand{\arraystretch}{1.25}
  \begin{tabularx}{0.96\textwidth}{@{}l l >{\RaggedRight\arraybackslash}X@{}}
    \toprule
    \textbf{Méthode} & \textbf{Ressource} & \textbf{Description} \\
    \midrule
    \endfirsthead
    \toprule
    \textbf{Méthode} & \textbf{Ressource} & \textbf{Description} \\
    \midrule
    \endhead
    POST & \texttt{/auth/login} & Authentification → retour JWT \\
    POST & \texttt{/auth/register} & Inscription citoyenne \\
    GET & \texttt{/auth/me} & Profil utilisateur courant \\
    GET & \texttt{/incidents} & Liste paginée et filtrée \\
    POST & \texttt{/incidents} & Création multipart (JSON + photo) \\
    GET & \texttt{/incidents/map} & Données légères pour Leaflet \\
    GET & \texttt{/incidents/recent} & Les 10 plus récents \\
    PATCH & \texttt{/incidents/\{id\}/status} & Mise à jour statut (ADMIN) \\
    GET & \texttt{/stats/dashboard} & Agrégats KPI \\
    GET & \texttt{/alerts} & Liste alertes (ADMIN) \\
    \bottomrule
  \end{tabularx}
\end{table}

\chapter{Réalisation — Frontend et intégration}
\label{chap:frontend}

\section{Architecture Next.js~14}

\subsection{App Router et pages}
\begin{itemize}
  \item \texttt{app/page.tsx} — Landing page, présentation ;
  \item \texttt{app/dashboard/page.tsx} — Tableau de bord (ADMIN / CITIZEN) ;
  \item \texttt{app/incidents/page.tsx} — Listing incidents, filtres ;
  \item \texttt{app/incidents/new/page.tsx} — Formulaire de création ;
  \item \texttt{app/map/page.tsx} — Carte Leaflet interactive ;
  \item \texttt{app/auth/page.tsx} — Authentification (login/register).
\end{itemize}

\subsection{Client HTTP et intercepteurs}
\begin{lstlisting}[basicstyle=\ttfamily\footnotesize,language=TypeScript,caption={Exemple : lib/api.ts avec intercepteur JWT.},label=lst:api_ts]
import axios from 'axios';

const api = axios.create({
  baseURL: process.env.NEXT_PUBLIC_API_URL ||
           'http://localhost:8080/api/v1',
});

api.interceptors.request.use((config) => {
  const token = localStorage.getItem('urbanops_token');
  if (token) {
    config.headers.Authorization = `Bearer ${token}`;
  }
  return config;
});

export default api;
\end{lstlisting}

\subsection{Composants React}
\begin{itemize}
  \item \texttt{components/IncidentForm.tsx} — Formulaire création ;
  \item \texttt{components/IncidentMap.tsx} — Intégration Leaflet ;
  \item \texttt{components/Dashboard.tsx} — Graphiques Recharts ;
  \item \texttt{components/AlertBanner.tsx} — Notification alertes (rouge HIGH, vert RESOLVED).
\end{itemize}

\urbanfig{fig:landing}{figures/landing_page_screenshot.png}{0.9}{Capture d'écran — page d'accueil (landing).}
\urbanfig{fig:dashboard}{figures/dashboard_screenshot.png}{0.92}{Capture d'écran — tableau de bord avec KPI et statuts.}
\urbanfig{fig:map}{figures/map_interactive_screenshot.png}{0.9}{Capture d'écran — carte Leaflet interactive des incidents actuels.}

\chapter{Sécurité, IA et middlewares}
\label{chap:api}

\section{Sécurité — JWT et BCrypt}

\subsection{Filtre d'authentification}
\begin{lstlisting}[style=javaurban,caption={JwtAuthenticationFilter — valide JWT à chaque requête.},label=lst:jwt_filter]
@Component
public class JwtAuthenticationFilter extends OncePerRequestFilter {
  @Override
  protected void doFilterInternal(HttpServletRequest request,
      HttpServletResponse response, FilterChain filterChain)
      throws ServletException, IOException {
    String token = extractToken(request);
    if (token != null && jwtProvider.validateToken(token)) {
      String username = jwtProvider.getUsernameFromJWT(token);
      UserDetails user = userDetailsService.loadUserByUsername(username);
      UsernamePasswordAuthenticationToken auth =
        new UsernamePasswordAuthenticationToken(user, null,
                                              user.getAuthorities());
      SecurityContextHolder.getContext().setAuthentication(auth);
    }
    filterChain.doFilter(request, response);
  }
}
\end{lstlisting}

\subsection{Chiffrement mots de passe}
Les mots de passe sont chiffrés lors de l'inscription via \texttt{BCryptPasswordEncoder} avec un nombre de rounds configurable (ex. 12).

\section{Module IA — Gemini + Repli déterministe}

\subsection{Flux de classification}
\begin{enumerate}
  \item Réception texte de l'incident (titre + description) ;
  \item Appel \textbf{Google Gemini API} avec prompts JSON structurés ;
  \item \textbf{Parsing JSON} → extraction sévérité, autorité recommandée ;
  \item \textbf{Fallback} : si timeout ou quota dépassé, classement par mots-clés prédéfinis.
\end{enumerate}

\subsection{Gestion des erreurs}
\begin{itemize}
  \item Timeout : repli immédiat (défaut \texttt{LOW}) ;
  \item Quota : retentative avec backoff exponentiel ;
  \item Parse error : log et assignation manuelle.
\end{itemize}

\section{Les quatre middlewares}

\subsection{1. REST (HTTP/JSON) — Principal}
\begin{itemize}
  \item Expose tous les endpoints (\texttt{/api/v1/*}) ;
  \item Stateless, scalable horizontalement ;
  \item Sécurisé par JWT ;
  \item Documentation OpenAPI via Swagger.
\end{itemize}

\subsection{2. JMS (ActiveMQ) — Messagerie asynchrone}
\begin{itemize}
  \item \textbf{Avantage} : découplage AlertService → EmailService ;
  \item Queues : \texttt{alerts-queue} pour mails en attente ;
  \item Retry automatique, persistance.
\end{itemize}

\subsection{3. RMI — Appels distants Java}
\begin{itemize}
  \item Service \texttt{AIAnalysisService} exporté en RMI ;
  \item Port : 1099 ;
  \item URL : \texttt{rmi://localhost:1099/AIAnalysisService} ;
  \item Permet intégration systèmes Java distants.
\end{itemize}

\subsection{4. SOAP — Services web XML}
\begin{itemize}
  \item WSDL public : \texttt{/api/v1/soap/urbanops.wsdl} ;
  \item Services : IncidentService SOAP, AlertService SOAP ;
  \item Compatibilité systèmes legacy (données XML).
\end{itemize}

\chapter{Tests, couverture et analyse SonarQube}
\label{chap:tests}

\section{Stratégie JUnit~5 et Mockito}

\subsection{Types de tests}
\begin{itemize}
  \item \textbf{Unitaires} : tests de services isolés, mocks de dépôts ;
  \item \textbf{Intégration} : \texttt{@SpringBootTest}, context applicatif complet ;
  \item \textbf{Repository} : requêtes JPA, projections ;
  \item \textbf{API} : tests HTTP via \texttt{MockMvc} ou TestRestTemplate.
\end{itemize}

\subsection{Exemple test unitaire}
\begin{lstlisting}[style=javaurban,caption={Test JUnit5 — IncidentService.createIncident().},label=lst:junit5]
@ExtendWith(MockitoExtension.class)
class IncidentServiceTest {

  @Mock
  private IncidentRepository incidentRepository;

  @Mock
  private AIAnalysisService aiService;

  @InjectMocks
  private IncidentService incidentService;

  @Test
  void testCreateIncident_Success() {
    // Arrange
    IncidentRequest request = new IncidentRequest("Nid", "Desc");
    when(aiService.analyze("Desc")).thenReturn(
      new AIResponse("HIGH", "Voirie"));
    when(incidentRepository.save(any())).thenAnswer(
      invocation -> invocation.getArgument(0));

    // Act
    IncidentResponse response =
      incidentService.createIncident(request);

    // Assert
    Assertions.assertNotNull(response.getId());
    Assertions.assertEquals("HIGH", response.getSeverity());
  }
}
\end{lstlisting}

\section{JaCoCo — Couverture de code}

\subsection{Configuration pom.xml}
\begin{lstlisting}[language=XML,caption={Fragment pom.xml — configuration JaCoCo.},label=lst:jacoco]
<plugin>
  <groupId>org.jacoco</groupId>
  <artifactId>jacoco-maven-plugin</artifactId>
  <version>0.8.8</version>
  <executions>
    <execution>
      <goals><goal>prepare-agent</goal></goals>
    </execution>
    <execution>
      <id>report</id>
      <phase>test</phase>
      <goals><goal>report</goal></goals>
    </execution>
    <execution>
      <id>jacoco-check</id>
      <goals><goal>check</goal></goals>
      <configuration>
        <rules>
          <rule>
            <element>BUNDLE</element>
            <excludes><exclude>*Test</exclude></excludes>
            <limits>
              <limit><counter>LINE</counter>
                <value>COVEREDRATIO</value>
                <minimum>0.70</minimum>
              </limit>
            </limits>
          </rule>
        </rules>
      </configuration>
    </execution>
  </executions>
</plugin>
\end{lstlisting}

\subsection{Génération rapport}
\begin{lstlisting}[style=bashurban,caption={Génération rapport JaCoCo.},label=lst:jacoco_cmd]
cd backend
mvn clean verify
# Rapport HTML sous: target/site/jacoco/index.html
\end{lstlisting}

\urbanfig{fig:jacoco}{figures/jacoco_report_screenshot.png}{0.95}{Capture — rapport JaCoCo global (couverture par package et classe).}

\section{SonarQube — Analyse statique qualité}

\subsection{Configuration}
\begin{itemize}
  \item \textbf{Serveur SonarQube} : \texttt{http://localhost:9000} (local) ou cloud ;
  \item \textbf{Authentification} : \texttt{sonar.login} token ;
  \item \textbf{Projet clé} : \texttt{ma.urbanops:backend}.
\end{itemize}

\subsection{Lancement analyse}
\begin{lstlisting}[style=bashurban,caption={Analyse SonarQube via Maven.},label=lst:sonar_cmd]
cd backend
mvn clean verify sonar:sonar \
  -Dsonar.projectKey=ma.urbanops:backend \
  -Dsonar.host.url=http://localhost:9000 \
  -Dsonar.login=YOUR_TOKEN
\end{lstlisting}

\subsection{Indicateurs mesurés}
\begin{table}[H]
  \centering
  \caption{Grille d'interprétation — indicateurs SonarQube.}
  \label{tab:sonar}
  \renewcommand{\arraystretch}{1.2}
  \begin{tabularx}{0.96\textwidth}{@{}l l X@{}}
    \toprule
    \textbf{Indicateur} & \textbf{Cible} & \textbf{Description} \\
    \midrule
    Fiabilité (Bugs) & \color{SuccessGreen}A & Zéro bug critique ou bloquant \\
    Vulnérabilités & \color{SuccessGreen}A & Pas d'injection SQL, XSS, etc. \\
    Hotspots sécurité & \color{SuccessGreen} Faible & Points critiques identifiés \\
    Couverture (JaCoCo) & \geq 70\% & Lignes couvertes par tests \\
    Duplications & < 3\% & Code dupliqué minimal \\
    Dépendances & À jour & Packages vulnérables éliminés \\
    \bottomrule
  \end{tabularx}
\end{table}

\urbanfig{fig:sonarqube}{figures/sonarqube_dashboard_screenshot.png}{0.95}{Capture — tableau de bord SonarQube (vue projet global).}

\section{Limitations et suites}
\begin{itemize}
  \item Tests réseau (@Disabled pour requêtes externes Gemini) ;
  \item Coverage des contrôleurs via \texttt{@SpringBootTest} le plus complet ;
  \item Test de charge (JMeter) non abordé mais recommandé.
\end{itemize}

\chapter{Déploiement et exploitation}
\label{chap:deploy}

\section{Lancement backend — développement}
\begin{lstlisting}[style=bashurban,caption={Démarrage du backend Spring Boot.},label=lst:backend_run]
cd backend

# Option 1 : Maven
mvn spring-boot:run

# Option 2 : Java direct
java -jar target/urbanops-backend-1.0.jar

# Vérification : swagger à http://localhost:8080/api/v1/swagger-ui.html
\end{lstlisting}

\section{Lancement frontend — développement}
\begin{lstlisting}[style=bashurban,caption={Démarrage du frontend Next.js.},label=lst:frontend_run]
cd frontend

npm install
npm run dev
# Accès : http://localhost:3000
\end{lstlisting}

\section{Variables d'environnement (.env)}
\begin{itemize}
  \item Backend : database URL, JWT secret, clé Gemini API, SMTP config ;
  \item Frontend : \texttt{NEXT\_PUBLIC\_API\_URL}, tokens d'analyse.
\end{itemize}

\section{Actuator — Healthchecks}
\begin{lstlisting}[basicstyle=\ttfamily\footnotesize,breaklines=true]
GET http://localhost:8080/api/v1/actuator/health

Réponse (exemple) :
{
  "status": "UP",
  "components": {
    "db": {"status": "UP"},
    "livenessState": {"status": "UP"},
    "readinessState": {"status": "UP"}
  }
}
\end{lstlisting}

\chapter{Conclusion et perspectives}
\label{chap:concl}

\section{Synthèse des apports}
UrbanOps démontre une \textbf{architecture robuste et maintenable} couplant technologies modernes (Spring~Boot~3, Next.js~14, PostgreSQL, Gemini). L'intégration de \textbf{quatre middlewares} (REST, JMS, RMI, SOAP) valide la polyvalence d'une plateforme urbaine. Les \textbf{tests unitaires JUnit~5}, \textbf{couverture JaCoCo} et \textbf{analyse SonarQube} garantissent une \textbf{qualité de production}.

Des succès clés :
\begin{itemize}
  \item Cycle de vie incident complèt (OPEN → RESOLVED) ;
  \item Classification IA automatisée avec secours déterministe ;
  \item Dashboard interactif et carte temps réel ;
  \item Sécurité JWT et encryption mots de passe.
\end{itemize}

\section{Difficultés rencontrées}
\begin{itemize}
  \item Gestion des quotas Gemini API et timeout réseau ;
  \item Synchronisation upload fichiers multipart et BDD ;
  \item Optimisation requêtes PostgreSQL (indices nécessaires) ;
  \item Intégration des quattro middlewares en parallèle.
\end{itemize}

\section{Perspectives et évolutions futures}
\begin{itemize}
  \item \textbf{Mobile} : Application React Native ou Flutter pour signalement sur terrain ;
  \item \textbf{Notifications} : Server-Sent Events (SSE) ou WebSockets push temps réel ;
  \item \textbf{IoT} : Capteurs IoT (humidité, température) intégrés aux incidents ;
  \item \textbf{Machine Learning} : Amélioration classification via fine-tuning modèle ;
  \item \textbf{Containers} : Déploiement Docker, orchestration Kubernetes ;
  \item \textbf{Analytics} : Tableaux analytiques avancés (ElasticSearch, Kibana) ;
  \item \textbf{Audit} : Logging blockchain pour immuabilité trace.
\end{itemize}

% =============================================================================
% BIBLIOGRAPHIE
% =============================================================================
\cleardoublepage
\chapter*{Bibliographie et webographie}
\addcontentsline{toc}{chapter}{Bibliographie et webographie}
\begin{thebibliography}{99}
  \bibitem{spring} Spring Framework Documentation — \url{https://spring.io/projects/spring-framework}
  \bibitem{springboot} Spring Boot 3 — \url{https://spring.io/projects/spring-boot}
  \bibitem{jpa} Jakarta Persistence API — \url{https://jakarta.ee/specifications/persistence/}
  \bibitem{nextjs} Next.js 14 Documentation — \url{https://nextjs.org/docs}
  \bibitem{typescript} TypeScript — \url{https://www.typescriptlang.org/docs/}
  \bibitem{tailwind} Tailwind CSS — \url{https://tailwindcss.com/docs}
  \bibitem{leaflet} Leaflet.js — \url{https://leafletjs.com/}
  \bibitem{postgresql} PostgreSQL — \url{https://www.postgresql.org/docs/}
  \bibitem{gemini} Google Generative AI — \url{https://ai.google.dev/docs}
  \bibitem{jwt} JWT.io — \url{https://jwt.io/}
  \bibitem{junit5} JUnit 5 User Guide — \url{https://junit.org/junit5/docs/current/user-guide/}
  \bibitem{mockito} Mockito — \url{https://javadoc.io/doc/org.mockito/mockito-core/latest/org/mockito/Mockito.html}
  \bibitem{jacoco} JaCoCo — \url{https://www.jacoco.org/}
  \bibitem{sonar} SonarQube Documentation — \url{https://docs.sonarsource.com/sonarqube/latest/}
  \bibitem{maven} Apache Maven — \url{https://maven.apache.org/}
  \bibitem{plantuml} PlantUML — \url{https://plantuml.com/}
  \bibitem{swagger} Swagger / OpenAPI — \url{https://swagger.io/}
\end{thebibliography}

% =============================================================================
% ANNEXES
% =============================================================================
\appendix

\chapter{Endpoints REST complets}
\label{app:endpoints}
\begin{longtable}{@{}l l >{\RaggedRight\arraybackslash}p{7.5cm}@{}}
  \toprule
  \textbf{Méthode} & \textbf{Ressource} & \textbf{Description} \\
  \midrule
  \endfirsthead
  \toprule
  \textbf{Méthode} & \textbf{Ressource} & \textbf{Description} \\
  \midrule
  \endhead
  POST & \texttt{/auth/login} & Authentification (email + password) \\
  POST & \texttt{/auth/register} & Nouvelle inscription citoyenne \\
  GET & \texttt{/auth/me} & Profil utilisateur courant (JWT requis) \\
  GET & \texttt{/incidents} & Liste paginée avec filtres (status, severity, date range) \\
  POST & \texttt{/incidents} & Créer incident (multipart: JSON + photo) \\
  GET & \texttt{/incidents/\{id\}} & Détail incident \\
  PATCH & \texttt{/incidents/\{id\}/status} & Modifier statut (ADMIN) \\
  GET & \texttt{/incidents/map} & Projection légère pour carte \\
  GET & \texttt{/incidents/recent} & Les 10 plus récents \\
  GET & \texttt{/incidents/stats} & Statistiques incidents \\
  GET & \texttt{/stats/dashboard} & KPI tableau de bord collectif \\
  GET & \texttt{/alerts} & Liste alertes existantes (ADMIN) \\
  POST & \texttt{/alerts} & Créer alerte manuelle \\
  GET & \texttt{/users} & Liste utilisateurs (ADMIN) \\
  DELETE & \texttt{/users/\{id\}} & Supprimer utilisateur (ADMIN) \\
  GET & \texttt{/categories} & Listing catégories \\
  GET & \texttt{/sectors} & Listing secteurs \\
  GET & \texttt{/actuator/health} & Healthcheck système \\
  \bottomrule
\end{longtable}

\chapter{Exemple configuration (anonymisé)}
\label{app:config}

\section{Backend — application.properties snippet}
\begin{lstlisting}[basicstyle=\ttfamily\footnotesize,language=bash,caption={Configuration backend essntielle.},label=lst:config_props]
server.port=8080
server.servlet.context-path=/api/v1

spring.datasource.url=jdbc:postgresql://localhost:5432/urbanops
spring.datasource.username=urbanops_user
spring.datasource.password=ENCRYPTED_PASSWORD

spring.jpa.hibernate.ddl-auto=update
spring.jpa.properties.hibernate.dialect=org.hibernate.dialect.PostgreSQLDialect

spring.mail.host=smtp.gmail.com
spring.mail.port=587
spring.mail.username=your-email@gmail.com
spring.mail.password=app-specific-password

app.jwt.secret=your-secret-key-min-256-bits-long
app.jwt.expiration=86400000

gemini.api.key=your-gemini-api-key
gemini.model=gemini-pro

file.upload-dir=./uploads
\end{lstlisting}

\section{Frontend — .env.local example}
\begin{lstlisting}[basicstyle=\ttfamily\footnotesize,language=bash,caption={Frontend Next.js — variables d'environnement.},label=lst:frontenv]
NEXT_PUBLIC_API_URL=http://localhost:8080/api/v1
\end{lstlisting}

\chapter{Schémas et diagrammes — guide export PlantUML}
\label{app:plantuml}

\section{Document PlantUML recommandé}
Exporter les fichiers PlantUML \texttt{.puml} versles formats PNG/PDF et les placer dans \texttt{docs/latex/figures/} avec les noms suggestifs :
\begin{itemize}
  \item \texttt{gantt\_planning.png} — Diagramme GANTT (phases et jalons) ;
  \item \texttt{pert\_taches.png} — Diagramme PERT (dépendances) ;
  \item \texttt{uml\_classes\_metier.png} — Classes métier ;
  \item \texttt{uml\_classes\_technique.png} — Classes techniques (Controller, Service, Repo) ;
  \item \texttt{uml\_cas\_usage.png} — Cas d'utilisation ;
  \item \texttt{uml\_composants.png} — Components diagram ;
  \item \texttt{uml\_sequence\_incident.png} — Sequence : création incident.
\end{itemize}

\chapter{Captures d'écran requises}
\label{app:captures}
Après implémentation, générer les captures d'écran :
\begin{itemize}
  \item \texttt{figures/landing\_page\_screenshot.png} — Page d'accueil ;
  \item \texttt{figures/dashboard\_screenshot.png} — Tableau de bord ;
  \item \texttt{figures/map\_interactive\_screenshot.png} — Carte interactive ;
  \item \texttt{figures/incident\_form\_screenshot.png} — Formulaire création ;
  \item \texttt{figures/jacoco\_report\_screenshot.png} — Rapport JaCoCo coverage ;
  \item \texttt{figures/sonarqube\_dashboard\_screenshot.png} — SonarQube réussite.
\end{itemize}

Les noms énumérés ici permettront au rapport LaTeX de les inclure automatiquement lors de la compilation.

\end{document}
